% Full proof annex for Durrett Section 3.2 and its required premises.
\section{Weak convergence of probability laws}
\label{sec:annex-p15-weak-convergence-of-laws}

The quantile, ramp-function, mapping, portmanteau and selection
arguments below retain their separate source routes
\cite[Section~3.2, pp.~116--125]{DurrettPTEAuthor2019}.

\begin{annexitem}{annex:p15.weak-quantile-foundations}{thm:p15-weak-quantile-foundations}{Countable jumps, inverse conventions and their measurable laws}
\begin{proof}
A bounded nondecreasing function has a left limit at every
point, equal to the supremum of its values strictly to the left.
For a right-continuous such function, a discontinuity is exactly
a positive jump \(F(x)-F(x-)\). For any finite ordered list
of points, the sum of their jumps is at most the total range,
which here is at most one: place left approximating points
before each listed point and telescope the monotone increments,
then pass to their limits. Thus at most \(m\) jumps can be
at least \(1/m\). Every positive jump belongs to one of these
finite sets, proving countability. Every nonempty interval
contains a point outside a countable set, so the continuity
points are dense.

Suppose the same proper CDF sequence converges weakly to \(F\)
and \(H\). At their common continuity points the limits agree.
Those points form a dense set, since the union of the two
discontinuity sets is countable. Given \(x\), choose common
continuity points decreasing to \(x\); right continuity gives
\(F(x)=H(x)\). The same argument also compares any two
right-continuous nondecreasing Helly limits, without endpoint
normalization, when both are limits of the same sequence.

For a proper \(F\) and \(0<u<1\), the two tail limits provide
real \(l<r\) with \(F(l)<u<F(r)\). The set defining
\(a=Q_F(u)\) is nonempty and bounded above, and the set
defining \(b=b_F(u)\) is nonempty and bounded below.
Both endpoints are finite and \(a\leq b\).
For the last inequality, every point with \(F(y)<u\) precedes
every point with \(F(z)>u\), by monotonicity; take the
supremum and infimum of these two sets.
If \(y<a\), the definition of the supremum supplies \(z>y\)
with \(F(z)<u\), whence \(F(y)<u\).
If \(y>a\), then \(F(y)\geq u\), since otherwise \(y\)
would belong to the defining set. Right continuity, applied
to \(a+1/j\), gives \(F(a)\geq u\).
It follows that \(a=\inf\{y:F(y)\geq u\}\).

For the exact ray identity, if \(u\leq F(x)\), monotonicity
excludes every \(y\geq x\) from \(\{F(y)<u\}\), so \(Q_F(u)\leq x\).
If \(u>F(x)\), right continuity supplies \(\eta>0\) with
\(F(x+\eta)<u\), giving \(Q_F(u)\geq x+\eta>x\).
This proves
\(\{Q_F\leq x\}=(0,1)\cap(0,F(x)]\).
All these inverse rays are Borel; because the rays generate the
Borel sigma-algebra of \(\mathbb R\), \(Q_F\) is measurable.
Their Lebesgue probabilities are \(F(x)\), so its pushforward
probability law has CDF \(F\). Uniqueness follows from the
already proved CDF law-determination theorem,
Theorem~\ref{thm:p15-cdf-laws}.
This supplies the full inverse-realization premise used in
source Theorem~1.2.2, with the source's strict inverse convention.
\end{proof}
\end{annexitem}

\begin{annexitem}{annex:p15.weak-quantile-coupling}{thm:p15-weak-quantile-coupling}{The source exceptional set and both limiting inequalities}
\begin{proof}
Write \(F=F_\infty\), \(a(u)=a_F(u)\), \(b(u)=b_F(u)\).
If \(0<u<v<1\), then \(b(u)\leq a(v)\).
Indeed any \(y<b(u)\) has \(F(y)\leq u<v\), so \(y\leq a(v)\);
let \(y\uparrow b(u)\). Hence the nonempty intervals
\((a(u),b(u))\) for distinct \(u\) are disjoint.
Choose in each one the first rational from a fixed enumeration
that it contains. These rationals are distinct. This injects
\(E=\{u:a(u)<b(u)\}\) into a countable set.
It is therefore countable, Borel and Lebesgue null.

If \(u\notin E\), then \(a(u)=b(u)=Q_F(u)\).
The inverse lemma gives \(F(y)<u\) for every \(y<Q_F(u)\).
For \(z>Q_F(u)=b(u)\), the infimum definition supplies some
\(w<z\) with \(F(w)>u\); monotonicity gives \(F(z)>u\).
Choose a continuity point \(y<Q_F(u)\).
The assumed CDF convergence implies \(F_n(y)<u\) eventually,
and hence \(Q_{F_n}(u)\geq y\).
By density of continuity points they can be chosen arbitrarily
close to \(Q_F(u)\) from below. Thus
\(\liminf_nQ_{F_n}(u)\geq Q_F(u)\).
For a continuity point \(z>Q_F(u)\), similarly \(F_n(z)>u\)
eventually, which implies \(Q_{F_n}(u)\leq z\).
Let these points approach from above to obtain
\(\limsup_nQ_{F_n}(u)\leq Q_F(u)\).

The inverse lemma gives every \(Y_n=Q_{F_n}\) its required
law on the single specified uniform probability space.
The pointwise convergence outside \(E\) proves the almost-sure
claim. This event is measurable also directly as the convergence
event of real measurable functions. Only the marginal laws
were prescribed; no claim about the original joint law follows.
\end{proof}
\end{annexitem}

\begin{annexitem}{annex:p15.weak-bounded-continuous}{thm:p15-weak-bounded-continuous}{The source coupling proof and continuous ramps}
\begin{proof}
Suppose first that the CDFs converge weakly, and use the
quantile coupling. For bounded continuous \(g\),
\(g(Y_n)\to g(Y_\infty)\) almost surely, and these functions
are bounded by the same finite constant.
Bounded convergence gives
\[
 \int g\,d\mu_n=\mathbb E g(Y_n)
    \longrightarrow\mathbb E g(Y_\infty)=\int g\,d\mu_\infty.
\]
Expectation depends only on the law by the established
integration-against-a-law theorem, so this also proves the
random-variable formulation for arbitrary original spaces.

Conversely, suppose all bounded continuous test integrals
converge. The displayed ramps are continuous, take values
in \([0,1]\), and obey
\[
 \mathbf1_{(-\infty,x]}\leq g_{x,\varepsilon}
       \leq\mathbf1_{(-\infty,x+\varepsilon]},\qquad
 \mathbf1_{(-\infty,x-\varepsilon]}
       \leq g_{x-\varepsilon,\varepsilon}
       \leq\mathbf1_{(-\infty,x]}.
\]
Consequently
\[
 \limsup_nF_n(x)\leq \int g_{x,\varepsilon}\,d\mu_\infty
                                      \leq F_\infty(x+\varepsilon),
\]
and
\[
 \liminf_nF_n(x)\geq\int g_{x-\varepsilon,\varepsilon}\,d\mu_\infty
                                      \geq F_\infty(x-\varepsilon).
\]
Let \(\varepsilon\downarrow0\). Right continuity makes the
first upper bound \(F_\infty(x)\) at every \(x\); the lower
bound becomes \(F_\infty(x-)\). At each continuity point
these values agree, proving exactly the stated CDF convergence.
\end{proof}
\end{annexitem}

\begin{annexitem}{annex:p15.weak-continuous-mapping}{thm:p15-weak-continuous-mapping}{Borel discontinuity set and the source composition argument}
\begin{proof}
We first prove the required measurability fact, even for an
arbitrary finite real function \(g\). For a nonempty open interval
\(I\), let its \(g\)-diameter be
\(\sup\{|g(y)-g(z)|:y,z\in I\}\), possibly \(+\infty\).
For each integer \(m\geq1\), let \(O_m\) be the union of
all open intervals having \(g\)-diameter less than \(1/m\).
This is open. If \(g\) is continuous at \(x\), a sufficiently
small interval around \(x\) has that diameter, so \(x\in O_m\)
for every \(m\). Conversely, if \(x\in\bigcap_mO_m\), choose,
for any requested tolerance, an interval containing \(x\)
whose diameter is smaller than that tolerance. Values of \(g\)
in a sufficiently small neighborhood of \(x\) then differ
from \(g(x)\) by less than the tolerance.
Thus the continuity set is \(\bigcap_mO_m\), and
\(D_g=\bigcup_mO_m^c\) is a countable union of closed sets.
This is the real-line discontinuity-set premise referred to
in source Exercise~1.3.6.

Its envelope proof can also be supplied directly. For
\(\delta>0\), define the extended functions
\[
 g^\delta(x)=\sup_{|y-x|<\delta}g(y),\qquad
 g_\delta(x)=\inf_{|y-x|<\delta}g(y).
\]
For every real \(a\), the set \(\{x:g^\delta(x)>a\}\)
is the union of the intervals \((y-\delta,y+\delta)\)
over points with \(g(y)>a\); it is open.
Similarly \(\{x:g_\delta(x)<a\}\) is open.
Thus both envelope functions are Borel measurable, including
possible infinite values. Set
\(g^0=\inf_{m\geq1}g^{1/m}\) and
\(g_0=\sup_{m\geq1}g_{1/m}\); these countable limits are
Borel and satisfy \(g_0(x)\leq g(x)\leq g^0(x)\).
Continuity at \(x\) is equivalent to equality of the two
envelopes there: equality forces their common value to be
the finite \(g(x)\), and the defining shrinking neighborhoods
then give the epsilon definition of continuity.
Conversely continuity bounds both envelopes by that value.
The set of strict inequality is Borel, for example by taking
the union of \(\{g_0<r<s<g^0\}\) over rational \(r<s\).
This retains the upper/lower-envelope route of the cited
exercise as well as the closed-set description above.

Now assume \(g\) is Borel and \(\mu_\infty(D_g)=0\).
Under the quantile coupling, intersect the probability-one
convergence event with \(\{Y_\infty\notin D_g\}\).
On this measurable event \(g(Y_n)\to g(Y_\infty)\).
For every bounded continuous \(f:\mathbb R\to\mathbb R\),
continuity of \(f\) implies \(D_{f\circ g}\subseteq D_g\);
in particular \(f(g(Y_n))\to f(g(Y_\infty))\) almost surely.
Bounded convergence and equality of the marginal laws give
\[
 \mathbb E f(g(X_n))\longrightarrow
                      \mathbb E f(g(X_\infty)).
\]
All compositions are measurable because \(g\) is Borel.
The bounded-continuous characterization, now applied to their
real laws, proves \(g(X_n)\Rightarrow g(X_\infty)\).
If \(g\) itself is bounded, apply bounded convergence directly
to \(g(Y_n)\) for the expectation conclusion. No finiteness
of expectations of an unbounded \(g\) was used or inferred.
\end{proof}
\end{annexitem}

\begin{annexitem}{annex:p15.weak-portmanteau}{thm:p15-weak-portmanteau}{Every implication in the source set characterization}
\begin{proof}
Assume (i), and realize the laws by the quantile coupling.
For an open \(G\), whenever \(Y_\infty\in G\) and
\(Y_n\to Y_\infty\), the sequence eventually lies in \(G\).
Thus
\(\liminf_n\mathbf1_G(Y_n)\geq\mathbf1_G(Y_\infty)\)
almost surely. Fatou's lemma gives
\[
 \liminf_n\mu_n(G)=\liminf_n\mathbb E\mathbf1_G(Y_n)
       \geq\mathbb E\mathbf1_G(Y_\infty)=\mu_\infty(G).
\]
This proves (ii). Taking complements, and using the identity
\(\mu_n(A^c)=1-\mu_n(A)\), shows that (ii) and (iii)
are equivalent, since a set is open exactly when its
complement is closed.

Assume (ii) and (iii), and let \(A\) be a Borel continuity set.
Set \(K=\overline A\), \(G=A^\circ\).
These are closed and open, respectively, and
\(G\subseteq A\subseteq K\), \(K\setminus G=\partial A\).
Because that boundary is \(\mu_\infty\)-null,
\(\mu_\infty(G)=\mu_\infty(A)=\mu_\infty(K)\).
The source's two comparisons are therefore
\[
 \liminf_n\mu_n(A)\geq\liminf_n\mu_n(G)
                   \geq\mu_\infty(G)=\mu_\infty(A)
\]
and
\[
 \limsup_n\mu_n(A)\leq\limsup_n\mu_n(K)
                   \leq\mu_\infty(K)=\mu_\infty(A).
\]
This proves (iv). Finally, if \(x\) is a continuity point
of \(F_\infty\), then
\(\mu_\infty(\{x\})=F_\infty(x)-F_\infty(x-)=0\).
The boundary of \((-\infty,x]\) is \(\{x\}\).
Applying (iv) to this ray proves \(F_n(x)\to F_\infty(x)\),
which is (i). The displayed implications form a complete
cycle and establish equivalence of all four conditions.
\end{proof}
\end{annexitem}

\begin{annexitem}{annex:p15.helly-selection}{thm:p15-helly-selection}{The source rational diagonal and the remaining finite measure}
\begin{proof}
Enumerate \(\mathbb Q\) as \(q_1,q_2,\ldots\).
Every sequence in \([0,1]\) has a convergent subsequence:
bisect \([0,1]\), retain a closed half containing infinitely
many terms, and repeat. The nested intervals have lengths
tending to zero and a common point, and increasing indices
selected from them give convergence to that point.
This supplies the scalar selection premise directly.

Starting with the full increasing index sequence, successively
take an increasing subsequence \(m_j(i)\) of the preceding
sequence such that \(F_{m_j(i)}(q_j)\) converges as \(i\to\infty\).
Set \(n_k=m_k(k)\). These indices strictly increase: nesting
gives \(m_{k+1}(k+1)\geq m_k(k+1)>m_k(k)\).
For each fixed \(j\), the tail of this diagonal is a
subsequence of the \(j\)-th selected sequence. Hence
\(F_{n_k}(q)\to G(q)\) for every rational \(q\).
The rational function \(G\) takes values in \([0,1]\) and
is nondecreasing, by passage to limits in the corresponding
inequalities for \(F_{n_k}\).

Define \(F(x)=\inf_{q>x,\ q\in\mathbb Q}G(q)\).
Its values belong to \([0,1]\), and it is nondecreasing.
If \(x_j\downarrow x\) with \(x_j>x\), the rational sets
\(\{q:q>x_j\}\) increase to \(\{q:q>x\}\).
Their infima consequently decrease to the infimum over their
union: any member of the union occurs at a finite stage.
Thus \(F(x_j)\to F(x)\). This proves right continuity;
the same conclusion is immediate if some \(x_j=x\).

Let \(x\) be a continuity point of \(F\), and fix \(\eta>0\).
Choose rationals \(r_1<r_2<x<s\) sufficiently close to \(x\)
that
\[
 F(x)-\eta<F(r_1)\leq F(r_2)\leq F(x)\leq F(s)<F(x)+\eta.
\]
Such choices follow from the continuity at \(x\) and density
of the rationals. By definition \(F(r_1)\leq G(r_2)\),
whereas monotonicity of \(G\) gives \(G(s)\leq F(s)\).
Therefore the rational convergences imply, eventually,
\[
 F(x)-\eta<F_{n_k}(r_2)\leq F_{n_k}(x)
                       \leq F_{n_k}(s)<F(x)+\eta.
\]
This establishes all continuity-point limits in the source
selection theorem.

The limits \(\beta,\gamma\) exist by bounded monotonicity.
If \(d=\gamma-\beta>0\), then
\(H(x)=(F(x)-\beta)/d\) is a proper CDF. The inverse
realization lemma constructs its probability law \(\eta\).
The finite measure \(\nu=d\eta\) has CDF \(F-\beta\)
and total mass \(d\). If \(d=0\), monotonicity forces \(F\)
to be constant, and take \(\nu=0\).
Uniqueness for positive \(d\) follows by normalization and
CDF uniqueness; for \(d=0\), a nonnegative measure of total
mass zero is zero. Subtracting the CDF limits at continuity
points \(r<s\) proves the interval formula.
The masses absent from the normalized endpoints are exactly
\(\beta\) at the negative end and \(1-\gamma\) at the positive
end, as quantified by these endpoint limits.
\end{proof}
\end{annexitem}

\begin{annexitem}{annex:p15.weak-tightness-selection}{thm:p15-weak-tightness-selection}{Both source directions and the uniform-tail equivalence}
\begin{proof}
Write \(H_n(M)=1-F_n(M)+F_n(-M)\).
Suppose first the source tightness condition holds, and
\(F_{n_k}\Rightarrow_v F\).
For \(\eta>0\), choose \(M\) with \(\limsup_nH_n(M)\leq\eta\).
Choose continuity points \(r<-M\) and \(s>M\) of \(F\).
Then
\[
 1-F(s)+F(r)
   =\lim_k\{1-F_{n_k}(s)+F_{n_k}(r)\}
   \leq\limsup_nH_n(M)\leq\eta.
\]
Let \(r\to-\infty\) and \(s\to\infty\) through continuity
points, which exist arbitrarily far in each direction.
This gives \(1-\gamma+\beta\leq\eta\).
Since both nonnegative summands have sum at most every
\(\eta>0\), \(\beta=0,\gamma=1\); the limit is proper.

Conversely suppose the source tightness condition fails.
Its logical negation supplies \(\eta_0>0\) with
\(\limsup_nH_n(M)>\eta_0\) for every finite \(M>0\).
Choose increasing indices \(n_k\) such that \(H_{n_k}(k)\geq\eta_0\).
Helly selection gives a further subsequence, indexed by
\(k_j\to\infty\), with \(F_{n_{k_j}}\Rightarrow_v F\).
For continuity points \(r<0<s\), eventually \(k_j>\max(-r,s)\),
so
\[
 1-F(s)+F(r)
 =\lim_j\{1-F_{n_{k_j}}(s)+F_{n_{k_j}}(r)\}
 \geq\liminf_jH_{n_{k_j}}(k_j)\geq\eta_0.
\]
Letting these points approach the two infinities gives
\(1-\gamma+\beta\geq\eta_0\). This selected limit is not proper.
The two implications prove the source equivalence.

For the uniform formulation, given \(\eta>0\), first choose
\(M\) with \(\limsup_nH_n(M)\leq\eta/2\).
There is \(N\) with \(H_n(M)\leq\eta\) for \(n\geq N\).
Choose \(R>M\) sufficiently large that
\(\mu_n(\{|x|\geq R\})\leq\eta\) for each of the finitely
many \(n<N\); this follows from continuity from above of
each probability measure and finiteness of its real variable.
For \(n\geq N\), the event \(\{|x|\geq R\}\) is contained
in \((-\infty,-M]\cup(M,\infty)\), proving the same bound.
Conversely \(H_n(M)\leq\mu_n(\{|x|\geq M\})\), so the
uniform condition implies the source one.

If every Helly limit is proper, applying Helly to any
subsequence supplies a proper weakly convergent further
subsequence. Conversely, a defective Helly subsequence
cannot have a further subsequence with a proper weak limit:
the two right-continuous limits would agree on their dense
common continuity set and hence everywhere, by the uniqueness
argument already proved. This proves the additional equivalence.
Finally a weakly convergent sequence has only its proper
limit as a possible Helly subsequential limit, by the same
uniqueness argument. It is therefore tight.
\end{proof}
\end{annexitem}

\begin{annexitem}{annex:p15.weak-coercive-tightness}{thm:p15-weak-coercive-tightness}{The source moment bound with explicit measurability}
\begin{proof}
For any \(L>0\), choose \(M\) from coercivity.
Pointwise \(L\mathbf1_{\{|x|\geq M\}}\leq\varphi(x)\).
Both functions are measurable and nonnegative, so integration
against each \(\mu_n\) gives
\[
 L\mu_n(\{|x|\geq M\})\leq\int\varphi\,d\mu_n\leq C.
\]
Divide by the finite positive \(L\), take the supremum over
\(n\), and then choose \(L\) large for any required tolerance.
The uniform formulation proves tightness. This also handles
\(C=0\) and possible infinite values of \(\varphi\);
no division by an infinite infimum is required.
The function \(\varphi(x)=|x|^p\) is Borel and coercive for
every \(p>0\), proving the moment consequence.
\end{proof}
\end{annexitem}

\begin{annexitem}{annex:p15.weak-metric-premises}{thm:p15-weak-metric-premises}{The complete Levy and Ky Fan constructions}
\begin{proof}
Call \(\varepsilon\geq0\) admissible for \((F,G)\) when the
two defining CDF inequalities hold. The value one is
admissible, since all CDF values lie in \([0,1]\).
Admissibility is preserved on increasing \(\varepsilon\),
by monotonicity of \(F\). Consequently every
\(\varepsilon>\rho(F,G)\) is admissible.
If \(\varepsilon\) is admissible for \((F,G)\), replacing
\(x\) by \(x-\varepsilon\) in the upper bound and by
\(x+\varepsilon\) in the lower bound gives
\[
 G(x-\varepsilon)-\varepsilon\leq F(x)
                      \leq G(x+\varepsilon)+\varepsilon.
\]
Thus the same values are admissible for \((G,F)\);
\(\rho\) is symmetric.
If \(\varepsilon,\delta\) are admissible for \((F,G),(G,H)\),
substitution in the two bounds gives
\[
 F(x-\varepsilon-\delta)-\varepsilon-\delta
 \leq H(x)\leq
 F(x+\varepsilon+\delta)+\varepsilon+\delta.
\]
Taking \(\varepsilon\downarrow\rho(F,G)\) and
\(\delta\downarrow\rho(G,H)\) proves the triangle inequality.
If \(\rho(F,G)=0\), every positive \(\varepsilon\) is
admissible. In particular \(G(x)\leq F(x+\varepsilon)+\varepsilon\);
right continuity gives \(G(x)\leq F(x)\).
Symmetry gives the reverse inequality, so \(F=G\).
The converse is immediate with \(\varepsilon=0\).
This proves all metric properties and the bound by one.

If \(\rho(F_n,F)\to0\), then for each fixed \(\eta>0\),
the value \(\eta\) is eventually admissible for \((F,F_n)\).
Hence
\[
 F(x-\eta)-\eta\leq\liminf_nF_n(x)
       \leq\limsup_nF_n(x)\leq F(x+\eta)+\eta.
\]
At continuity points of \(F\), let \(\eta\downarrow0\)
to obtain weak convergence.
Conversely suppose \(F_n\Rightarrow F\), and fix \(0<\eta<1\).
Choose continuity points \(a<b\) of \(F\) with
\(F(a)<\eta/2\) and \(F(b)>1-\eta/2\).
Insert finitely many continuity points to form a partition
\(a=x_0<x_1<\cdots<x_m=b\) whose mesh is less than \(\eta\).
Such a partition exists by density of continuity points,
for example by perturbing a sufficiently fine finite uniform
grid by less than a quarter of its mesh.
For all sufficiently large \(n\),
\(|F_n(x_j)-F(x_j)|<\eta/2\) at every partition point.
If \(x_j\leq x\leq x_{j+1}\), monotonicity gives
\[
 F(x-\eta)-\eta
 \leq F(x_j)-\eta
 \leq F_n(x)
 \leq F(x_{j+1})+\eta
 \leq F(x+\eta)+\eta.
\]
For \(x<a\), the lower bound is at most zero, while
\(F_n(x)\leq F_n(a)<\eta\), proving the upper bound.
For \(x>b\), \(F_n(x)\geq F_n(b)>1-\eta\) proves the
lower bound, and \(F(x+\eta)+\eta>1\) proves the upper.
Thus \(\eta\) is admissible for \((F,F_n)\) eventually;
\(\rho(F_n,F)\leq\eta\), giving the required metric limit.
The partition uses spatial mesh; it imposes no continuity
assumption at the jumps of \(F\).

For \(\alpha\), the value one is admissible and the
admissible set is upward closed. Let \(a=\alpha(X,Y)\).
Every \(a+1/j\) is admissible. The events
\(\{|X-Y|>a+1/j\}\) increase to \(\{|X-Y|>a\}\).
Continuity from below gives
\(\mathbb P(|X-Y|>a)\leq a\); thus the infimum is attained.
It follows that \(\alpha=0\) exactly when \(X=Y\)
almost surely. The expression is unchanged on almost-sure
equivalence classes and is symmetric.
If \(\varepsilon,\delta\) are admissible for \((X,Y),(Y,Z)\),
the triangle inequality for real numbers gives
\[
 \{|X-Z|>\varepsilon+\delta\}
 \subseteq\{|X-Y|>\varepsilon\}\cup\{|Y-Z|>\delta\}.
\]
The union bound shows \(\alpha(X,Z)\leq\varepsilon+\delta\).
Use the attained infima, or let the parameters decrease
to them, to prove the metric triangle inequality.

If \(\alpha(X_n,X)=a_n\to0\), then for a fixed \(\eta>0\)
and all large \(n\) with \(a_n<\eta\),
\(\mathbb P(|X_n-X|>\eta)\leq
\mathbb P(|X_n-X|>a_n)\leq a_n\to0\).
Conversely convergence in probability gives
\(\mathbb P(|X_n-X|>\eta)\leq\eta\) eventually for every
\(\eta>0\), so \(\alpha(X_n,X)\leq\eta\).

With \(a=\alpha(X,Y)\) and the exact attained bound,
the event \(\{X\leq x-a\}\), except for
\(\{|X-Y|>a\}\), is contained in \(\{Y\leq x\}\).
Conversely \(\{Y\leq x\}\), except for that same error
event, is contained in \(\{X\leq x+a\}\).
Therefore
\[
 F(x-a)-a\leq G(x)\leq F(x+a)+a
 \quad\hbox{for every real }x,
\]
and \(\rho(F,G)\leq a\).
The convergence-in-probability implication follows from
the two metric characterizations.
For the constant converse, the closed set
\(K=\{x:|x-c|\geq\eta\}\), \(\eta>0\), has zero mass
under \(\delta_c\). The closed-set portmanteau bound gives
\(\limsup_n\mathbb P(|X_n-c|\geq\eta)\leq0\),
and hence the desired strict-threshold convergence as well.
\end{proof}
\end{annexitem}

\begin{annexitem}{annex:p15.weak-subsequence-criterion}{thm:p15-weak-subsequence-criterion}{The source metric route and the direct CDF route}
\begin{proof}
For the source's metric route, suppose
\(\rho(F_n,F_\infty)\) did not tend to zero. Some positive
\(\eta\) and subsequence would have distance at least
\(\eta\) at every index. No further subsequence could
converge to \(F_\infty\) in that metric, and therefore
none could converge weakly to it, by the metric theorem.
This contradicts the hypothesis and establishes the source
argument from its now-proved premises.

There is also a direct CDF proof.
If \(F_n\) did not converge weakly to \(F_\infty\), there
would be a continuity point \(x\) of \(F_\infty\), a number
\(\eta>0\), and a subsequence \(n_k\) with
\(|F_{n_k}(x)-F_\infty(x)|\geq\eta\) for all \(k\).
No further subsequence of that one could converge weakly to
\(F_\infty\), contradicting the hypothesis.
This proves the same assertion independently of the metric route.

Under tightness, every subsequence has a further proper
weak limit by the selection theorem. If all such limits
have the target law, the first part applies. This gives
the stated practical criterion.
\end{proof}
\end{annexitem}

\begin{annexitem}{annex:p15.scheffe-density-limit}{thm:p15-scheffe-density-limit}{The source positive-part identity and exact event supremum}
\begin{proof}
All densities are integrable and finite almost everywhere.
On the union of their countably many exceptional null sets
they may be assigned the value zero; this changes none of the
laws, integrals or almost-everywhere convergence statements.
Let \(h_n=f_n-f_\infty\). It is integrable and has integral
zero. Its positive and negative parts therefore have equal
finite integrals, giving
\[
 \int|h_n|\,d\lambda
      =2\int h_n^-\,d\lambda
      =2\int(f_\infty-f_n)^+\,d\lambda.
\]
The last integrand tends to zero almost everywhere and is
bounded above by the integrable \(f_\infty\).
Dominated convergence proves the limit. This needs no
finite-total-mass assumption on \(\lambda\).

For any measurable \(A\), the integral of \(h_n\) over \(A\)
lies between \(-\int h_n^-d\lambda\) and
\(\int h_n^+d\lambda\). Hence its absolute value is at most
\(\tfrac12\int|h_n|d\lambda\). The measurable set
\(\{h_n\geq0\}\) attains the upper bound.
Taking the supremum proves the exact total-variation identity
in the source's convention. On the real line, apply the
event bound to every ray \((-\infty,x]\) to get convergence
of CDFs at every real \(x\), and thus weak convergence.
\end{proof}
\end{annexitem}

\begin{annexitem}{annex:p15.weak-elementary-asymptotics}{thm:p15-weak-elementary-asymptotics}{The logarithmic remainder and an independent central coefficient calculation}
\begin{proof}
For \(|t|\leq1/2\), the integral representation of the
logarithm gives
\[
 |\log(1+t)-t|
 =\left|\int_0^t\frac{-v}{1+v}\,dv\right|
 \leq t^2.
\]
Eventually every array entry lies in this range, so all
factors are positive and
\[
 \left|\sum_j\log(1+c_{n,j})-\sum_jc_{n,j}\right|
 \leq\sum_j|c_{n,j}|^2
 \leq\max_j|c_{n,j}|\sum_j|c_{n,j}|\longrightarrow0.
\]
Exponentiation proves the product limit.
The scalar source premise also follows directly: if
\(c_n\to0\), \(a_n\to\infty\), \(a_nc_n\to\ell\), then
\[
 |a_n\log(1+c_n)-a_nc_n|\leq |a_nc_n|\,|c_n|\to0,
\]
where \(a_n>0\) eventually. Thus \((1+c_n)^{a_n}\to e^\ell\),
including noninteger real exponents with their eventually
positive bases.

For the central coefficient put
\(I_n=\int_{-\pi/2}^{\pi/2}\cos^{2n}t\,dt\).
Integration by parts with factors
\(\cos^{2n-1}t\) and \(\cos t\,dt\) has zero endpoint term
and gives
\[
 I_n=(2n-1)(I_{n-1}-I_n),\qquad
 I_n=\frac{2n-1}{2n}I_{n-1}\quad(n\geq1).
\]
Since \(I_0=\pi\), multiplication of these finite identities
gives \(I_n=\pi(2n)!/(4^n(n!)^2)=\pi C_n\).

For \(|t|<\pi/2\),
\(\log\cos t=-\int_0^{|t|}\tan v\,dv\).
The derivative of \(\tan v\) is \(1/\cos^2v\geq1\), so
\(\tan v\geq v\) for \(v\geq0\) in this interval.
Thus \(\log\cos t\leq-t^2/2\).
Also \(\tan v/v\to1\) at zero, by the derivative at zero;
bounding this ratio above and below in the integral proves
\(\log\cos t/t^2\to-1/2\).

Define \(g_n(u)=\cos^{2n}(u/\sqrt n)\) on
\(|u|<\pi\sqrt n/2\) and zero outside.
These nonnegative measurable functions satisfy
\(g_n(u)\leq e^{-u^2}\), and for every fixed \(u\),
\(g_n(u)\to e^{-u^2}\).
The affine substitution \(u=\sqrt n\,t\) gives
\(\sqrt n I_n=\int_{\mathbb R}g_n(u)\,du\).
The already established Gaussian integral is
\(\int_{\mathbb R}e^{-u^2}du=\sqrt\pi\).
Dominated convergence therefore yields
\(\sqrt n I_n\to\sqrt\pi\), or
\(C_n\sqrt{\pi n}\to1\).
Only the one-variable affine substitution and Gaussian
normalization already proved in the manuscript are used.
\end{proof}
\end{annexitem}

\begin{annexitem}{annex:p15.weak-rademacher-example}{thm:p15-weak-rademacher-example}{Local binomial masses, Riemann sums and both parities}
\begin{proof}
For \(N=2m\), counting which \(m+r\) signs are positive gives
\[
 p_m(r)=\mathbb P(S_{2m}=2r)
       =4^{-m}\binom{2m}{m+r}\quad(-m\leq r\leq m).
\]
Its central value is \(p_m(0)=C_m\).
For \(0\leq r\leq m\), cancellation of factorials gives
\[
 \frac{p_m(r)}{C_m}
   =\prod_{j=1}^r\frac{m-j+1}{m+j}.
\]
Fix \(0<A<\infty\). For \(r\leq A\sqrt m\) and large \(m\),
all \((j-1)/m,j/m\) are at most \(1/2\).
The logarithmic remainder bound of the preceding lemma yields
\[
 \log\frac{p_m(r)}{C_m}
 =-\frac1m\sum_{j=1}^r(2j-1)+R_{m,r}
 =-\frac{r^2}{m}+R_{m,r},\qquad
 |R_{m,r}|\leq\frac2{m^2}\sum_{j=1}^rj^2.
\]
The last sum is at most \(r(r+1)^2\), so its upper bound
tends to zero uniformly for \(0\leq r\leq A\sqrt m\).
Symmetry gives the same statement for negative \(r\).
Combining this with \(C_m\sim1/\sqrt{\pi m}\), and setting
\(h_m=2/\sqrt{2m}\), proves
\[
 \sup_{\substack{r\in\mathbb Z\\ |2r/\sqrt{2m}|\leq A}}
 \left|\frac{p_m(r)}{h_m}
             -\phi(2r/\sqrt{2m})\right|\longrightarrow0.
\]
The source local-mass premise needed by the example is thus
proved uniformly on each bounded scaled interval.

For fixed real \(a<b\), sum over the lattice
\(2\mathbb Z/\sqrt{2m}\) in \([a,b]\).
The contribution of the preceding uniform error is at most
\((b-a+2h_m)\) times that error, and tends to zero.
The remaining sum is the Riemann sum
\(\sum_{x\in[a,b]\cap h_m\mathbb Z}h_m\phi(x)\).
Uniform continuity of \(\phi\) on a slightly enlarged bounded
interval shows convergence to \(\int_a^b\phi(x)\,dx\):
on each lattice interval the difference from its left-endpoint
value is bounded by its modulus of continuity at \(h_m\),
and the two partial endpoint intervals have lengths at most
\(h_m\) and bounded integrand. This supplies all endpoint
and Riemann-sum details.

The signs have mean zero and variance one. Independence gives
\(\mathbb E S_N^2=N\), and Chebyshev gives
\(\mathbb P(|S_N|/\sqrt N>A)\leq A^{-2}\).
For any fixed \(x\), first restrict the CDF calculation to
\([-A,x]\), where \(A>|x|\), and use the bounded-interval
limit. The omitted left tail has upper limit at most \(A^{-2}\);
the corresponding Gaussian tail tends to zero as \(A\to\infty\).
This proves the CDF limit for even \(N\).

For odd indices, put \(T_m=S_{2m}/\sqrt{2m}\) and
\(d_m=\sqrt{2m/(2m+1)}\).
Since \(S_{2m+1}=S_{2m}+\varepsilon_{2m+1}\),
\[
 \mathbb P\!\left(T_m\leq x/d_m-1/\sqrt{2m}\right)
 \leq\mathbb P(S_{2m+1}/\sqrt{2m+1}\leq x)
 \leq\mathbb P\!\left(T_m\leq x/d_m+1/\sqrt{2m}\right).
\]
Both moving thresholds tend to \(x\).
For any \(\eta>0\) they eventually lie between \(x-\eta\)
and \(x+\eta\); monotonicity and the even CDF limits then
bound their lower and upper limits by \(\Phi(x-\eta)\)
and \(\Phi(x+\eta)\).
Letting \(\eta\downarrow0\), using continuity of \(\Phi\),
proves the odd CDF limit. The result now holds for all \(N\),
with no imported characteristic-function or CLT theorem.
\end{proof}
\end{annexitem}

\begin{annexitem}{annex:p15.weak-first-examples}{thm:p15-weak-first-examples}{Empirical laws, moving atoms, rare events and birthdays}
\begin{proof}
For the empirical example, the preceding complete
Glivenko--Cantelli theorem,
Theorem~\ref{thm:p15-glivenko-cantelli}, supplies one
measurable event of probability one on which
\(\sup_x|F_n(x)-F(x)|\to0\).
Each empirical CDF on each outcome is a proper CDF.
Hence on that event the CDFs converge at every real point,
in particular at all continuity points. This proves the
source example also for discontinuous \(F\).

For the translated variable,
\(\mathbb P(X+1/n\leq x)=F(x-1/n)\).
The arguments \(x-1/n\) increase to \(x\) from below, so
the limit is \(F(x-)\). It equals \(F(x)\) at continuity
points and is strictly smaller at every positive jump,
which accounts for every discontinuity of a CDF.

For the geometric variable, the event of never succeeding is
a measurable decreasing intersection with probability
\(\lim_m(1-p)^m=0\). The prescribed value one on that event
makes \(X_p\) finite everywhere and changes no probabilities.
For an integer \(m\geq1\), its event \(X_p\geq m\) agrees
up to this null set with failure of the first \(m-1\) trials.
Its probability is therefore \((1-p)^{m-1}\).
For \(x\geq0\), integer-valuedness gives
\[
 \mathbb P(pX_p>x)=(1-p)^{\lfloor x/p\rfloor}.
\]
For \(x>0\), its exponent tends to infinity and
\(-p\lfloor x/p\rfloor\to-x\).
The scalar logarithmic limit already proved gives \(e^{-x}\).
At \(x=0\) the probability equals one for every \(p\),
and at \(x<0\) the CDF is zero. These are the CDF limits
of the rate-one exponential probability law. Its density
\(e^{-x}\mathbf1_{\{x>0\}}\) integrates to one by direct
integration.

For birthdays, the first repeat occurs no later than \(N+1\),
by the pigeonhole principle. The event \(T_N>m\) is that the
first \(m\) observations are all distinct.
Given a distinct prefix of length \(j\leq N\), exactly \(j\)
of the \(N\) possible next values would be repeats.
Independence of the next draw gives the conditional
continuation probability \(1-j/N\); equivalently, counting
ordered distinct prefixes gives \(N(N-1)\cdots(N-m+1)/N^m\).
Multiplication proves the stated product for \(1\leq m\leq N+1\),
including its zero factor at \(m=N+1\). The other cases
follow directly from the event definition.

For fixed \(x>0\), put \(m_N=\lfloor x\sqrt N\rfloor\).
Eventually \(1\leq m_N\leq N\), and the product-limit
lemma applies to \(c_{N,j}=-j/N\), \(1\leq j<m_N\):
\[
 \max_{j<m_N}|c_{N,j}|\to0,\quad
 \sum_{j=1}^{m_N-1}c_{N,j}
       =-\frac{m_N(m_N-1)}{2N}\longrightarrow-x^2/2,
\]
and the sums of absolute values are bounded.
Therefore
\(\mathbb P(T_N/\sqrt N>x)=\mathbb P(T_N>m_N)\to e^{-x^2/2}\).
For \(x=0\) the probability is one, and negative \(x\)
have zero limiting CDF. The limiting function is proper
and continuous; equivalently its density on \(x>0\) is
\(xe^{-x^2/2}\), whose integral is one by its antiderivative
\(-e^{-x^2/2}\).
The finite product at \(N=365,m=22\) is
\(\prod_{j=1}^{21}(1-j/365)=0.5243046923\ldots\).
The source's leading approximation is
\(e^{-484/730}=0.5152960451\ldots\).
These calculations describe the specified independent uniform
model, without extending it to nonuniform birth frequencies.
\end{proof}
\end{annexitem}

\begin{annexitem}{annex:p15.weak-central-order-statistic}{thm:p15-weak-central-order-statistic}{The complete combinatorial density and source Scheffe route}
\begin{proof}
Sort the finite list of observations and choose its middle
entry. For any real \(x\), its event of being at most \(x\)
is
\[
 \{V_{n+1}\leq x\}
   =\left\{\sum_{i=1}^{2n+1}\mathbf1_{\{U_i\leq x\}}\geq n+1\right\},
\]
which is measurable. This proves measurability of the order
statistic even when values coincide.
There are no ties almost surely: for each pair, independence
and integration over one uniform coordinate give
\(\mathbb P(U_i=U_j)=\int_0^1\mathbb P(U_j=t)\,dt=0\).
There are only finitely many pairs.

Off that null set, partition the event that the median lies
in a Borel set \(B\subseteq(0,1)\) by its index \(i\) and
the subset \(J\) of exactly \(n\) other indices below it.
The remaining \(n\) indices must lie above it. For each fixed
choice, product independence and Tonelli give probability
\[
 \int_B t^n(1-t)^n\,dt.
\]
There are \(2n+1\) choices of \(i\) and \(\binom{2n}{n}\)
choices of \(J\), so summing the disjoint cases proves the
stated density. It integrates to one because these cases
exhaust the median event, and the median lies in \((0,1)\)
almost surely. This is the source index-counting proof
with its infinitesimal-density step expressed as an exact
integral over arbitrary Borel sets.

The already proved affine density formula, with
\(x=1/2+y/(2\sqrt{2n})\), gives, on \(|y|<\sqrt{2n}\),
\[
 \begin{aligned}
 f_{Y_n}(y)
 &=\frac{(2n+1)\binom{2n}{n}}{2\sqrt{2n}}
   \left(\frac12+\frac{y}{2\sqrt{2n}}\right)^n
   \left(\frac12-\frac{y}{2\sqrt{2n}}\right)^n\\
 &=C_n(1-y^2/(2n))^n
                      \frac{2n+1}{2n}\sqrt{\frac n2}.
 \end{aligned}
\]
The inverse affine image of the original support gives the
support restriction; set the density to zero at both endpoints
and outside. For each fixed \(y\), the support eventually
contains \(y\). The central-coefficient limit gives
\(C_n\sqrt{n/2}\to1/\sqrt{2\pi}\), the scalar product
limit gives \((1-y^2/(2n))^n\to e^{-y^2/2}\), and
\((2n+1)/(2n)\to1\).
The density therefore tends pointwise to \(\phi(y)\).
Every density integrates to one, and the limiting Gaussian
density is normalized. Scheffe's theorem applies and gives
total-variation convergence, hence the asserted weak limit.
This retains the source affine-density and Scheffe argument,
with the central-coefficient premise independently established.
\end{proof}
\end{annexitem}

\begin{annexitem}{annex:p15.weak-convergence-boundaries}{thm:p15-weak-convergence-boundaries}{Explicit countermodels and both directions of escaping mass}
\begin{proof}
For the point masses at \(1/n\), each continuity point of
the CDF of \(\delta_0\) is different from zero, and the CDFs
eventually agree there. Thus the weak convergence holds.
The event \(\{0\}\) has probability zero under every
\(\delta_{1/n}\) and one under \(\delta_0\), so the
total-variation distance is one, the largest possible value.
For \(G=(0,\infty)\), the probabilities are one and zero;
for its closed complement \(K\), they are zero and one.
These give the strict portmanteau inequalities in the
respective directions.
The function \(\mathbf1_{\{0\}}\) is continuous at every
nonzero point and discontinuous at zero. Its images in
this example are the distinct fixed laws \(\delta_0,\delta_1\),
so they do not converge to the asserted image of the limit.

For the fair sign, the law of \(-X\) equals that of \(X\),
giving weak convergence of the constant-in-law sequence.
But \(|-X-X|=2\) surely, proving failure of convergence in
probability on the given space.

For the variable with masses \(1/n\) at \(n\) and \(1-1/n\)
at zero, at a fixed negative \(x\) its CDF is zero.
For a fixed positive \(x\), it is \(1-1/n\) for all
sufficiently large \(n\). Hence its law converges weakly
to \(\delta_0\); its mean is exactly \(n(1/n)=1\).
The continuous unbounded test \(g(x)=x\) therefore fails
to have convergent expectations to that of the limit.

Each \(f_n=\mathbf1_{[n,n+1]}\) has Lebesgue integral one.
At each fixed real point it is eventually zero, so its
pointwise limit integrates to zero. Indeed its mass on
every fixed bounded interval is eventually zero, and
its CDF at each fixed \(x\) is eventually zero. This is
a defective Helly limit and gives no proper weak limit.

Finally the displayed mixture \(F_n\) is the CDF of
\(a\delta_n+b\delta_{-n}+c\eta\), where \(\eta\) has
CDF \(G\). For fixed \(x\), eventually \(x<n\) and
\(x\geq-n\), giving \(F_n(x)=b+cG(x)\).
Its negative and positive endpoint limits are \(b\)
and \(b+c\), respectively, and its finite remaining
measure is \(c\eta\).
If \(a+b>0\), this particular Helly limit is not proper,
so the tightness theorem rules out tightness.
If \(a=b=0\), the sequence is the single fixed probability
law \(\eta\), which is tight by its two CDF tail limits.
This proves the exact mass-loss statement in both directions.
\end{proof}
\end{annexitem}
